% XeLaTeX 编译；自编演示题，使用时替换为核对过的原卷内容。
\documentclass[11pt,a4paper]{ctexart}
\usepackage[margin=20mm]{geometry}
\usepackage{amsmath,amssymb,xcolor}
\definecolor{subjectgreen}{HTML}{087B52}
\setlength{\parindent}{0pt}
\setlength{\parskip}{0.6em}
\linespread{1.18}
\newcommand{\qlabel}[1]{\par\nopagebreak\textbf{\textcolor{subjectgreen}{#1}}\quad}
\begin{document}
{\LARGE\bfseries\textcolor{subjectgreen}{数学教师版 · 排版示例}}

自编演示题；使用时替换为核对过的原卷内容。

\section*{第 1 题}
已知函数 $f(x)=x^2-2x+2$，求该函数的最小值，并说明理由。

\qlabel{考点}二次函数的最值。

\qlabel{答案}最小值为 $1$。

\qlabel{解析}先配方，转化为顶点式：
\[
f(x)=x^2-2x+2=(x-1)^2+1.
\]
因为 $(x-1)^2\geq 0$，当且仅当 $x=1$ 时取等，所以函数的最小值为 $1$。

\section*{第 2 题}
已知 $a>0$，考虑函数 $g(x)=x+\frac{a}{x}$（$x>0$）。求其最小值，并说明取等条件。

\qlabel{考点}基本不等式与取等条件。

\qlabel{答案}最小值为 $2\sqrt a$。

\qlabel{解析}由于 $x>0$、$a>0$，两项均为正，可以使用基本不等式：
\[
g(x)=x+\frac{a}{x}\geq 2\sqrt{x\cdot\frac{a}{x}}=2\sqrt a.
\]
当 $x=\frac{a}{x}$，即 $x=\sqrt a$ 时取等。因此，最小值为 $2\sqrt a$。
\end{document}
